\section{The Executed-Cover Design Space}\label{sec:design}
\paragraph{Four axes}
We classify a cover policy along four axes (Table~\ref{tab:cover-taxonomy}). \emph{Termination} and \emph{state} determine what business work cover performs: a disposable copy, read-only access, an isolated clone or production state. \emph{Feedback visibility} determines whether cover responses enter later model requests and can change the real planner's decisions. The \emph{dispatcher}, the process that decides when cover is sent, determines whether cover inherits the timing of honest calls. Honest calls follow a model decision; a controller that inserts cover without waiting for the model creates a timing signature that the first three axes do not capture. Our live study varies the dispatcher while holding termination (an isolated clone) and feedback (hidden) fixed; an earlier state control varies termination and feedback.

\begin{table*}[t]\centering\small
\caption{Cover taxonomy. Each policy is classified by where it terminates and what state it touches, whether its responses reach the real planner, and which process dispatches it. The last column names the evaluation in which the policy is measured; S1 is the earlier state control and Section~\ref{sec:dispatch} the live timing study.}\label{tab:cover-taxonomy}
\setlength{\tabcolsep}{4pt}
\begin{tabularx}{\textwidth}{@{}P{.17\textwidth}P{.26\textwidth}P{.13\textwidth}P{.2\textwidth}Y@{}}\toprule
Policy & Terminates at / state touched & Feedback to planner & Dispatcher & Measured in \\\midrule
Transport cover & Tunnel endpoint before the service; no business state & No & Sender schedule & STP~\cite{apthorpe2019stp}; WF padding~\cite{juarez2016wtfpad,gong2020front} \\
Proxy discard & Defender proxy computes on a disposable copy & No & Controller & S1 \\
Read cover & Real service; reads only & No & Controller & S1 \\
Isolated execution & Real service on a persistent clone & No / yes & Controller & S1 (hidden and visible); all live cover arms \\
Live write & Real service on production state & No & Controller & S1 (stress control) \\
Research decoys & Real third-party Web destinations & Not located & Decoy generator & WRA~\cite{jeong2026wra} \\\midrule
Immediate & Isolated clone & No & Controller, directly after an honest response & Section~\ref{sec:dispatch} (calibration) \\
History-matched & Isolated clone & No & Controller, training-gap and duration samples; sequential or concurrent & Section~\ref{sec:dispatch} \\
Planner shadow & Isolated clone; separate planner context & No & Second planner on an independent ticket & Section~\ref{sec:dispatch} \\\bottomrule
\end{tabularx}\end{table*}


\paragraph{The execution contract}
Because executed cover performs work, a policy must satisfy a contract as well as a privacy target. The contract names the objects and transitions cover may touch, the quota, token and latency budgets it may consume, and the feedback it may expose to the real planner. Isolation prevents direct writes to production state but still consumes compute and can change scheduling; a response hidden from the planner does not undo the operation that produced it. A requested goal can complete while another order, a shared quota or a ledger is wrong, so the contract separates three outcomes: requested-goal completion, violations of declared business invariants, and resource use. For a synthetic fixture whose intended transition is fully specified, exact state equality is the invariant; in production, invariants must admit legitimate concurrent updates.

\paragraph{Where executed cover sits among defenses}
Name hiding and relays change destination association; padding changes lengths; pacing changes timing~\cite{wang2017walkietalkie,gong2020front,juarez2016wtfpad}. These layers remain useful, and Table~\ref{tab:cover-necessity} shows when transport cover suffices. Executed cover is the remaining option when the observed hop ends at a service the defender does not control. Model services add their own signals, such as token streams, cache behavior and token counts~\cite{weiss2024prompt,whisperleak2025,earlybird2025,timewilltell2024,carlini2024timing}; a tool-side cover policy does not shape later model requests. Appendix~\ref{app:defense} summarizes what each defense layer requires the defender to control, the costs to measure together and the exposure that remains.
